\section*{Chapitre \ref{ch:g5:fractions} --- Les fractions comme nombres}

\begin{solution}{exo:g5:fractions:1}
$\frac25$~: deux pas depuis $0$~; $\frac55$~: sur $1$~;
$\frac85$~: trois pas après $1$~; $\frac{10}{5}$~: sur $2$.
\end{solution}

\begin{solution}{exo:g5:fractions:2}
$\frac12 = \frac48$~; \quad $\frac34 = \frac68$~; \quad
$\frac25 = \frac{4}{10}$~; \quad $\frac69 = \frac23$.
\end{solution}

\begin{solution}{exo:g5:fractions:3}
$\frac82 = 4$~; \quad $\frac94$ n'est pas entier ($9$ n'est pas un
multiple de $4$)~; \quad $\frac{15}{3} = 5$~; \quad
$\frac{20}{5} = 4$.
\end{solution}

\begin{solution}{exo:g5:fractions:4}
$0.7$~; \quad $0.41$~; \quad $0.5$~; \quad $0.75$~; \quad $0.2$.
\end{solution}

\begin{solution}{exo:g5:fractions:5}
$0.9 = \frac{9}{10}$~; \quad $0.13 = \frac{13}{100}$~; \quad
$2.5 = \frac{25}{10}$~; \quad $0.75 = \frac{75}{100} = \frac34$.
\end{solution}

\begin{solution}{exo:g5:fractions:6}
$\frac47 < \frac67$ (même \omterm{def:g4:fractions:def}{dénominateur}). \quad
$\frac35 > \frac37$ (même \omterm{def:g4:fractions:def}{numérateur}~: les cinquièmes sont plus
grands). \quad
$\frac25 < \frac12 < \frac{7}{12}$, donc $\frac25 < \frac{7}{12}$
(repère $\frac12$). \quad
$\frac98 > 1 > \frac{11}{12}$ (repère $1$).
\end{solution}

\begin{solution}{exo:g5:fractions:7}
Décimaux~: $0.5$~; $0.4$~; $0.75$~; $0.6$~; $0.2$. Ordre~:
\[
\frac15 < 0.4 < \frac12 < 0.6 < \frac34 .
\]
\end{solution}

\begin{solution}{exo:g5:fractions:8}
Instrument~: $\frac14$ de $28 = 7$. Sport~: $\frac12$ de $28 = 14$.
Le groupe sport est plus grand.
\end{solution}

\begin{solution}{exo:g5:fractions:9}
$\frac14$~h $= 15$~min~; $\frac12$~h $= 30$~min~; $\frac34$~h $=
45$~min~; $\frac{1}{60}$~h $= 1$~min.
\end{solution}

\begin{solution}{exo:g5:fractions:10}
Tom~: $\frac35$ de $50 = 30$~cL. Ana~: $\frac12$ de $60 = 30$~cL.
Ils ont bu la même quantité~! Les \omterm{def:g4:fractions:def}{fractions} nues
($\frac35 > \frac12$) comparent des parts de bouteilles
\emph{différentes} --- seules les quantités effectives se
comparent.
\end{solution}

\begin{solution}{exo:g5:fractions:11}
Entre $\frac12$ et $1$~: par exemple $\frac34$ (c'est-à-dire
$0.75$). Entre $\frac12 = 0.5$ et $\frac34 = 0.75$~: par exemple
$0.6 = \frac35$.
\end{solution}
